\begin{flushright}
  \fechaclase{26 de agosto de 2026}
\end{flushright}

\paragraph{Matriz definida no negativa:}
Una matriz simétrica $P\in\mathbb{R}^{n\times n}$ se dice definida no negativa
si
\[
  x^{T}Px\geq 0
\]
para todo
\[
  x\in\mathbb{R}^{n}
\]

\paragraph{Lema de Bellman:}
Las siguientes afirmaciones son equivalentes

\begin{enumerate}
  \item $S$ es definida no negativa

  \item Se puede representar como
    \[
      S=HH^{T}
    \]
    para alguna matriz $H$

  \item Los valores propios de $S$ no son negativos
    \[
      \lambda_i(S)\geq 0, \qquad i=1,2,\ldots
    \]

  \item Existe una matriz simétrica
    \[
      M\in\mathbb{R}^{n\times n}
    \]
    tal que
    \[
      S=M^{2}
    \]
    $S$ es la raíz cuadrada de $M$ y se denota por
    \[
      S^{1/2}=M
    \]
\end{enumerate}

\paragraph{Matriz definida positiva:}
Una matriz simétrica es definida positiva si todos sus valores propios son
mayores que cero. Se denota como
\[
  S>0, \qquad S\in\mathbb{R}^{n\times n}
\]

\paragraph{Matriz semidefinida positiva:}
Es una matriz simétrica con valores propios mayores o iguales a cero. Se denota
como
\[
  S\geq 0, \qquad S\in\mathbb{R}^{n\times n}
\]

\paragraph{El cociente de Rayleigh:}
Se define como la función
\[
  R(x)\colon\mathbb{C}^{n}\longrightarrow\mathbb{R}
\]
con la matriz $A\in\mathbb{C}^{n\times n}$ y el vector
$x\in\mathbb{C}^{n}$:
\[
  R(x)=\frac{x^{T}Ax}{x^{T}x}
\]
para los valores propios de la matriz se cumple
\[
  \lambda_{\min}(A)<R(x)<\lambda_{\max}(A)
\]
Sustituyendo
\[
  \lambda_{\min}(A)
  <\frac{x^{T}Ax}{x^{T}x}
  <\lambda_{\max}(A)
\]
\[
  \lambda_{\min}(A)
  <\frac{x^{T}Ax}{\lVert x\rVert^{2}}
  <\lambda_{\max}(A)
\]
\[
  \lambda_{\min}(A)\lVert x\rVert^{2}
  <x^{T}Ax
  <\lambda_{\max}(A)\lVert x\rVert^{2}
\]

\subsubsection{El complemento de Schur}

Una división de matrices y submatrices se refiere a la partición de la matriz
como
\[
  A=
  \begin{bmatrix}
    A_{11} & A_{12}\\
    A_{21} & A_{22}
  \end{bmatrix}
  \qquad
  B=
  \begin{bmatrix}
    B_{11} & B_{12}\\
    B_{21} & B_{22}
  \end{bmatrix}
\]
donde las submatrices $A_{i,j}$ y $B_{i,j}$ tienen el mismo tamaño, entonces
\[
  A\pm B=
  \begin{bmatrix}
    A_{11}\pm B_{11} & A_{12}\pm B_{12}\\
    A_{21}\pm B_{21} & A_{22}\pm B_{22}
  \end{bmatrix}
\]
\[
  AB=
  \begin{bmatrix}
    A_{11}B_{11}+A_{12}B_{21} &
    A_{11}B_{12}+A_{12}B_{22}\\
    A_{21}B_{11}+A_{22}B_{21} &
    A_{21}B_{12}+A_{22}B_{22}
  \end{bmatrix}
\]

\paragraph{Teorema de Albert:}
Sea $S$ una matriz cuadrada particionada
\[
  S=
  \begin{bmatrix}
    S_{11} & S_{12}\\
    S_{21} & S_{22}
  \end{bmatrix}
\]
con $S_{11}\in\mathbb{R}^{n\times n}$ y
$S_{22}\in\mathbb{R}^{m\times m}$ entonces $S\geq 0$ si y solo si
\[
  S_{11}\geq 0
\]
\[
  S_{11}S_{11}^{+}S_{12}=S_{12}
\]
\[
  S_{22}-S_{12}^{T}S_{11}^{+}S_{12}\geq 0
\]